% page 432 of the facsimile (image p-431 of the scan)
% block 152.4 x 246.3 mm ; left margin 8.0 mm
\pgc{-12.700mm}{0.000mm}{
\l{}
\l{}
\l{\cel{3}424}
\l{}
\l{\sou{partenokarpio}. (\sou{biol.})\cel{2}Nomo quan Noll propozis en 1902}
\l{\cel{3}por indikar la modo di su-developo che ula frukti.}
\l{}
\l{\sou{partio}. Inter-amuzo, inter-distrakto da grupo de personi.}
\l{\cel{3}DEFS.}
\l{}
\l{\sou{partenogenezo}. Modo di genito ye qua femino reproduktas sua}
\l{\cel{3}speco sen fekundigesir da maskulo. - DEFIRS.}
\l{}
\l{\sou{partero}. I. Ta parto di gardeno ube la sulo esas tre plata}
\l{\cel{3}e glata, sen-arbora, e dividita per bedi. - II. (\sou{teatro})}
\l{\cel{3}Ta parto di spektaklo-chambrego qua esas inter la orkes-}
\l{\cel{3}treyo e la fundo di la teatro ube, olim, la spektanti}
\l{\cel{3}permanis stacanta. - DFIR.}
\l{}
\l{\sou{particiono}.\cel{2}(\sou{muziko).}\cel{2}Asembluro de omna parti di kanto}
\l{\cel{3}od orkestro. - DEFIRS.}
\l{}
\l{\sou{participo.}\cel{3}(\sou{gram.)}\cel{2}Vortospeco qua esas parte verba e par-}
\l{\cel{3}te adjektiva - modo ne-personala di la verbo, qua expre-}
\l{\cel{3}sas la ago forme di adjektivo. - DEFIS.}
\l{}
\l{\sou{partikulo.}\cel{2}I. Parto mikrega di korpo. - II.\cel{1}\sou{(gram.)}\cel{2}Mikra}
\l{\cel{3}(kurta) vorto, generale du-silaba, nevarianta, qua uzesas}
\l{\cel{3}por la kompozo di altra vorti. - III. La partikulo di}
\l{\cel{3}nobeleso : la prepoziciono qua esas situita avan (ante)}
\l{\cel{3}la nomo di multa familii nobela. - DEFIS.}
\l{}
\l{\sou{partikulara}.\cel{2}Qua esas la proprajo exkluziva di ulu, di ulo.}
\l{\cel{3}DEFIRS.}
\l{}
\l{\sou{partiso}. Grupo de personi qui adoptis (od adheris a) la}
\l{\cel{3}sama principi di konduto. - DEFIRS.}
\l{}
\l{\sou{partitivo.}\cel{2}(\sou{gram.}) Qua indikas parto ek toto. - DEFIS.}
\l{}
\l{\sou{partizano.}\cel{2}Olima shaft-armo, di qua la lamo esis plu larja}
\l{\cel{3}e plu akuta kam olta di la halbardo ordinara. - DEFIS.}
\l{}
\l{\sou{parturar}. (\sou{trans.)}\cel{2}(\sou{aludante patrino}). Emisar de su (fi-}
\l{\cel{3}lio konceptita e qua atingis la nasko-instanto). - EFIS.}
\l{}
\l{\sou{paruo}.\cel{1}\sou{(zool.}) Mikra ucelo, ek la klaso \textquotedbl{}paseri\textquotedbl{}, gracila,}
\l{\cel{3}di qua la bela plumaro esas bunta. - L.\cel{1}\sou{parus.}\cel{2}SL.}
\l{}
\l{\sou{pasar}. (\sou{netrans.)}\cel{2}I. Durar sua iro, sua fluo. - II. Venan-}
\l{\cel{3}te de... abutar a...\cel{2}- DEFIS.}
\l{}
\l{\sou{pasabla.}\cel{1}Qua esas aceptebla o tolerebla, quankam minima}
\l{\cel{3}qualesale o quantesale. - DEFIS.}
\l{}
\l{\sou{pasajar.}\cel{1}(\sou{netrans.})\cel{2}Voyajar per navo, pos lokacir sua pla-}
\l{\cel{3}so dure di la voyajo e pagir po lo. - DEFIRS.}
}
